\chapter{任务定义与系统构成}

\section{当前要完成的完整任务}

当前现场任务是：桌面上散放多个塑料瓶；左侧工作臂从桌面选择并抓起一个瓶子，右侧工作臂抓住瓶盖并将其旋下；左臂把瓶身投入瓶身回收盒，右臂把瓶盖投入瓶盖回收盒；随后两臂继续处理下一个瓶子。\ev{E01,E10}

\begin{table}[H]
\centering
\caption{单次任务的开始、完成、失败和中止条件}
\begin{tabularx}{\textwidth}{>{\bfseries}p{26mm}X p{39mm}}
\toprule
类别 & 当前定义 & 证据状态\\
\midrule
开始条件 & 双臂处于可运行初始姿态，桌面存在待处理瓶子，三路画面可用 & 运行配置可核验；初始摆放未形成固定测试协议\\
任务完成 & 瓶盖与瓶身分离，并分别进入两个回收盒 & 现场人工观察；没有自动判定\\
典型失败 & 未抓到瓶、倒着抓、无盖仍空拧、瓶盖未分离、投错盒或掉落 & 前两项有现场描述；频次未记录\\
中止条件 & 人员主动停止、设备异常或安全检查不通过 & 采集与控制软件具备停止/健康检查能力\\
\bottomrule
\end{tabularx}
\end{table}

\section{硬件构成与双臂分工}

\begin{table}[H]
\centering
\caption{系统硬件与职责}
\begin{tabularx}{\textwidth}{>{\bfseries}p{30mm}p{39mm}X}
\toprule
组成 & 主要职责 & 已核验事实与限制\\
\midrule
左侧工作臂与夹爪 & 选取、夹持、抬起瓶身；旋盖时稳定瓶身；投放瓶身 & 双臂输出中包含左侧 6 个关节和 1 个夹爪量；职责由任务和现场观察确认\\
右侧工作臂与夹爪 & 接近瓶口、夹住瓶盖、旋转与分离；投放瓶盖 & 含连续旋转动作；没有独立力觉/触觉记录\\
顶部相机 & 看清桌面全局、瓶子分布和两臂相对位置 & 正式训练和部署均使用；原始画面 640$\times$480\\
左腕相机 & 提供左夹爪、瓶身和接触近景 & 正式训练和部署均使用\\
右腕相机 & 提供瓶盖、右夹爪和旋转近景 & 注意力统计中平均占比最高\\
操作者示教臂 & 采集人员用双手直接演示动作 & 支持连续采集、重采、回位与多种触发方式\\
工作台与回收盒 & 提供随机瓶堆和瓶/盖分类投放区域 & 尺寸、坐标和固定摆放尚未形成公开标定表\\
\bottomrule
\end{tabularx}
\end{table}

\section{软硬件闭环}

\figfull[.98\textwidth]{software_data_pipeline.pdf}{第一年度形成的数据—训练—部署闭环。专用采集、数据编辑、版本发布、模型训练、现场运行和问题回流已经连通；当前软件能力不反推历史每条数据都使用过全部后来新增功能。}

采集人员通过专用软件遥操作 ALOHA，软件同时保存机器人状态、多视角视频和任务信息；数据进入编辑中心后进行浏览、裁剪、标签、版本生成和发布；训练平台读取固定版本数据，保存模型并记录训练过程；正式模型再部署到机器人。现场失败应当回流到数据分类与补采，而不是仅靠反复延长训练时间。

\section{控制周期、状态与动作}

正式系统以 50 Hz 发送控制指令，即每 20 毫秒一个控制时刻。机器人状态和动作均为 14 个量：
\[
\mathbf{s}_t,\mathbf{a}_t\in\mathbb{R}^{14}
=[q^L_{1:6},g^L,q^R_{1:6},g^R],
\]
其中 $q^L_{1:6}$ 和 $q^R_{1:6}$ 分别是左右工作臂 6 个关节位置，$g^L,g^R$ 是两只夹爪的位置。模型一次预测未来 $H=50$ 个时刻的动作片段，因此输出维度为 $50\times14$。\ev{E06,E07}

\plain{可以把它想成机器人每次先写好“接下来一秒怎么动”的小计划；执行到中间时，就开始准备下一张小计划，减少两张计划之间的停顿。}

运行时，系统在动作片段开始后第 25 个时刻启动后台重规划，并在满足切换条件后接入下一片段。该机制是“边执行边准备下一段”，不是把多次预测简单平均。关节动作会受到硬件范围限制；夹爪电流上限也有明确设置。\ev{E07}

\limitbox{当前代码没有自动判断“瓶盖是否脱离”“瓶身是否进盒”或“任务是否成功”。碰撞次数、瓶盖转角和接触力也未记录。因此成功条件仍依赖人工观察。}
